\documentclass[10pt,a4paper]{amsart}
\usepackage{amsmath,amssymb,amsthm,graphicx}
% Retained minimal preamble for the Stacks Project chapter more-morphisms.tex (More on Morphisms).
% Verbatim excerpt of the chapter preamble: Xy-pic setup (native preamble lines 14-18) and the native
% theorem declarations with the mathematical macros (native preamble lines 154-217). Copied mechanically;
% nothing added, removed or redefined. The wrapper class and the amsmath/amssymb/amsthm/graphicx base
% packages are supplied by the assembler, so the class/externaldocument part of the native preamble is omitted.
% For commutative diagrams we use Xy-pic
\usepackage[all]{xy}

% We use 2cell for 2-commutative diagrams.
\xyoption{2cell}
% Theorem environments.
%
\theoremstyle{plain}
\newtheorem{theorem}[subsection]{Theorem}
\newtheorem{proposition}[subsection]{Proposition}
\newtheorem{lemma}[subsection]{Lemma}

\theoremstyle{definition}
\newtheorem{definition}[subsection]{Definition}
\newtheorem{example}[subsection]{Example}
\newtheorem{exercise}[subsection]{Exercise}
\newtheorem{situation}[subsection]{Situation}

\theoremstyle{remark}
\newtheorem{remark}[subsection]{Remark}
\newtheorem{remarks}[subsection]{Remarks}

\numberwithin{equation}{subsection}

% Macros
%
\def\lim{\mathop{\mathrm{lim}}\nolimits}
\def\colim{\mathop{\mathrm{colim}}\nolimits}
\def\Spec{\mathop{\mathrm{Spec}}}
\def\Hom{\mathop{\mathrm{Hom}}\nolimits}
\def\Ext{\mathop{\mathrm{Ext}}\nolimits}
\def\SheafHom{\mathop{\mathcal{H}\!\mathit{om}}\nolimits}
\def\SheafExt{\mathop{\mathcal{E}\!\mathit{xt}}\nolimits}
\def\Sch{\mathit{Sch}}
\def\Mor{\mathop{\mathrm{Mor}}\nolimits}
\def\Ob{\mathop{\mathrm{Ob}}\nolimits}
\def\Sh{\mathop{\mathit{Sh}}\nolimits}
\def\NL{\mathop{N\!L}\nolimits}
\def\CH{\mathop{\mathrm{CH}}\nolimits}
\def\proetale{{pro\text{-}\acute{e}tale}}
\def\etale{{\acute{e}tale}}
\def\QCoh{\mathit{QCoh}}
\def\Ker{\mathop{\mathrm{Ker}}}
\def\Im{\mathop{\mathrm{Im}}}
\def\Coker{\mathop{\mathrm{Coker}}}
\def\Coim{\mathop{\mathrm{Coim}}}

% Boxtimes
%
\DeclareMathSymbol{\boxtimes}{\mathbin}{AMSa}{"02}

%
% Macros for moduli stacks/spaces
%
\def\QCohstack{\mathcal{QC}\!\mathit{oh}}
\def\Cohstack{\mathcal{C}\!\mathit{oh}}
\def\Spacesstack{\mathcal{S}\!\mathit{paces}}
\def\Quotfunctor{\mathrm{Quot}}
\def\Hilbfunctor{\mathrm{Hilb}}
\def\Curvesstack{\mathcal{C}\!\mathit{urves}}
\def\Polarizedstack{\mathcal{P}\!\mathit{olarized}}
\def\Complexesstack{\mathcal{C}\!\mathit{omplexes}}
% \Pic is the operator that assigns to X its picard group, usage \Pic(X)
% \Picardstack_{X/B} denotes the Picard stack of X over B
% \Picardfunctor_{X/B} denotes the Picard functor of X over B
\def\Pic{\mathop{\mathrm{Pic}}\nolimits}
\def\Picardstack{\mathcal{P}\!\mathit{ic}}
\def\Picardfunctor{\mathrm{Pic}}
\def\Deformationcategory{\mathcal{D}\!\mathit{ef}}
\usepackage{cleveref}
\crefname{originalreference}{original reference}{original references}
\Crefname{originalreference}{Original reference}{Original references}
\makeatletter
\newlabel{morphisms-lemma-composition-quasi-affine}{{01SN}{source}{}{sourcepointer-1}{}}
\newlabel{morphisms-lemma-composition-quasi-affine@cref}{{[originalreference][0][]01SN}{[1][0][]source}{}{}{}}
\makeatother
\crefname{originalreference}{original reference}{original references}
\Crefname{originalreference}{Original reference}{Original references}
\makeatletter
\newlabel{morphisms-lemma-image-proper-scheme-closed}{{01W6}{source}{}{sourcepointer-2}{}}
\newlabel{morphisms-lemma-image-proper-scheme-closed@cref}{{[originalreference][0][]01W6}{[1][0][]source}{}{}{}}
\makeatother
\crefname{originalreference}{original reference}{original references}
\Crefname{originalreference}{Original reference}{Original references}
\makeatletter
\newlabel{morphisms-lemma-integral-universally-closed}{{01WM}{source}{}{sourcepointer-3}{}}
\newlabel{morphisms-lemma-integral-universally-closed@cref}{{[originalreference][0][]01WM}{[1][0][]source}{}{}{}}
\makeatother
\crefname{originalreference}{original reference}{original references}
\Crefname{originalreference}{Original reference}{Original references}
\makeatletter
\newlabel{morphisms-lemma-quasi-compact-scheme-theoretic-image}{{01R8}{source}{}{sourcepointer-4}{}}
\newlabel{morphisms-lemma-quasi-compact-scheme-theoretic-image@cref}{{[originalreference][0][]01R8}{[1][0][]source}{}{}{}}
\makeatother
\crefname{originalreference}{original reference}{original references}
\Crefname{originalreference}{Original reference}{Original references}
\makeatletter
\newlabel{properties-lemma-quasi-affine}{{01P9}{source}{}{sourcepointer-5}{}}
\newlabel{properties-lemma-quasi-affine@cref}{{[originalreference][0][]01P9}{[1][0][]source}{}{}{}}
\makeatother
\title{The spectrum of global functions of an integral cover of a quasi-affine scheme}
\author{The Stacks Project Authors}
\date{stacks-project:more-morphisms.tex:local-snapshot-20261011; GNU FDL 1.2 or later}
\begin{document}
\maketitle
\noindent\textit{Editorial note.} Counted span is the author's own proof environment, source lines 13973-14010, immediately following the statement at 13960-13971.\par\medskip
\noindent\textit{Editorial note.} Five cross-chapter pointers (Morphisms, Properties) are recorded in assembly.external\_references.\par\medskip
\noindent\textit{Editorial note.} The cartesian square of the statement is displayed verbatim in the statement span.\par\medskip
\noindent\textit{Allowed standard prerequisites (editorial).}\begin{itemize}
\item Integral morphisms are universally closed. Morphisms, lemma morphisms-lemma-integral-universally-closed, cited in the proof. Provides closedness of the morphism, which is what makes the scheme-theoretic image argument work.
\item Scheme-theoretic image of a quasi-compact morphism. Morphisms, lemma morphisms-lemma-quasi-compact-scheme-theoretic-image, cited at proof lines 14007-14008. Identifies the image with the spectrum of global functions, giving density of U in it.
\item Openness of the structural arrows for quasi-affine schemes. Properties, lemma properties-lemma-quasi-affine, cited at proof lines 13997-13998. Shows the vertical arrows of the square are open, so that the comparison morphism j is an open immersion.
\item Composition of quasi-affine morphisms. Morphisms, lemma morphisms-lemma-composition-quasi-affine, cited in the proof. Gives quasi-affineness of U once the cartesian square is established.
\end{itemize}\medskip
\ifdefined\hypertarget\hypertarget{sourcepointer-1}{}\fi
\noindent\textit{Source reference.} Original reference 01SN: focused-seeds/geometry/S9-morph-integral-over-quasi-affine/ref-tags.txt, source lines 5--5. Navigation pointer only: the counted proof cites this statement by \textbackslash{}ref as a standing input; the deduction the counted proof has to make is carried out inside the counted span, and no part of that deduction is delegated to this pointer.\par\medskip
\ifdefined\hypertarget\hypertarget{sourcepointer-2}{}\fi
\noindent\textit{Source reference.} Original reference 01W6: focused-seeds/geometry/S9-morph-integral-over-quasi-affine/ref-tags.txt, source lines 6--6. Navigation pointer only: the counted proof cites this statement by \textbackslash{}ref as a standing input; the deduction the counted proof has to make is carried out inside the counted span, and no part of that deduction is delegated to this pointer.\par\medskip
\ifdefined\hypertarget\hypertarget{sourcepointer-3}{}\fi
\noindent\textit{Source reference.} Original reference 01WM: focused-seeds/geometry/S9-morph-integral-over-quasi-affine/ref-tags.txt, source lines 7--7. Navigation pointer only: the counted proof cites this statement by \textbackslash{}ref as a standing input; the deduction the counted proof has to make is carried out inside the counted span, and no part of that deduction is delegated to this pointer.\par\medskip
\ifdefined\hypertarget\hypertarget{sourcepointer-4}{}\fi
\noindent\textit{Source reference.} Original reference 01R8: focused-seeds/geometry/S9-morph-integral-over-quasi-affine/ref-tags.txt, source lines 8--8. Navigation pointer only: the counted proof cites this statement by \textbackslash{}ref as a standing input; the deduction the counted proof has to make is carried out inside the counted span, and no part of that deduction is delegated to this pointer.\par\medskip
\ifdefined\hypertarget\hypertarget{sourcepointer-5}{}\fi
\noindent\textit{Source reference.} Original reference 01P9: focused-seeds/geometry/S9-morph-integral-over-quasi-affine/ref-tags.txt, source lines 9--9. Navigation pointer only: the counted proof cites this statement by \textbackslash{}ref as a standing input; the deduction the counted proof has to make is carried out inside the counted span, and no part of that deduction is delegated to this pointer.\par\medskip
\begin{lemma}
\label{lemma-integral-over-quasi-affine}
Let $X$ be a quasi-affine scheme. Let $f : U \to X$ be an integral
morphism. Then $U$ is quasi-affine and the diagram
$$
\xymatrix{
U \ar[r] \ar[d] & \Spec(\Gamma(U, \mathcal{O}_U)) \ar[d] \\
X \ar[r] & \Spec(\Gamma(X, \mathcal{O}_X))
}
$$
is cartesian.
\end{lemma}

\begin{proof}
The scheme $U$ is quasi-affine because integral morphisms are affine,
affine morphisms are quasi-affine, a scheme is quasi-affine if and only if
the structure morphism to $\Spec(\mathbf{Z})$ is quasi-affine, and
compositions of quasi-affine morphisms are quasi-affine.
The first two statements follow immediately from the definition
and the third is
Morphisms, Lemma \ref{morphisms-lemma-composition-quasi-affine}.
Set $U' =
X \times_{\Spec(\Gamma(X, \mathcal{O}_X))} \Spec(\Gamma(U, \mathcal{O}_U))$
and consider the extended diagram
$$
\xymatrix{
U \ar[r]_j \ar[rd] & U' \ar[d] \ar[r] &
\Spec(\Gamma(U, \mathcal{O}_U)) \ar[d] \\
& X \ar[r] & \Spec(\Gamma(X, \mathcal{O}_X))
}
$$
The morphism $j$ is closed by
Morphisms, Lemma \ref{morphisms-lemma-image-proper-scheme-closed}
combined with the fact that an integral morphism is universally closed
(Morphisms, Lemma \ref{morphisms-lemma-integral-universally-closed})
and the fact that the vertical arrows are in the diagram are separated.
On the other hand, $j$ is open because the horizontal
arrows in the diagram of the lemma are open by
Properties, Lemma \ref{properties-lemma-quasi-affine}.
Thus $j$ identifies $U$ with an open and closed subscheme of $U'$.
If $U \not = U'$ then $U$ isn't dense in $U'$ and a fortiori
not dense in the spectrum of $\Gamma(U, \mathcal{O}_U)$.
However, the scheme theoretic image of
$U$ in $\Spec(\Gamma(U, \mathcal{O}_U))$ is $\Spec(\Gamma(U, \mathcal{O}_U))$
because any ideal in $\Gamma(U, \mathcal{O}_U)$
cutting out a closed subscheme through which $U$
factors would have to be zero.
Hence $U$ is dense in $\Spec(\Gamma(U, \mathcal{O}_U))$ for example by
Morphisms, Lemma \ref{morphisms-lemma-quasi-compact-scheme-theoretic-image}.
Thus $U = U'$ and we win.
\end{proof}

\end{document}
